\begin{frame}
    \frametitle{Introduction}
    \begin{columns}
        \column[t]{5cm}
    \begin{itemize}
    \item Multiple new reactor designs require \gls{HALEU} fuel, which allows for: 
    \begin{itemize}
        \item Longer cycle times
        \item Higher burnups 
    \end{itemize}
    \item<1-> To meet the \gls{HALEU} demand, the U.S. \gls{DOE} has proposed two methods \cite{griffith_overview_2020}:
    \begin{itemize}
        \item Recovery and downblending of \gls{HEU}
        \item Enrichment of natural uranium
    \end{itemize}
    \end{itemize}

    \column[t]{5cm}
    \begin{table}
        \centering
        \caption{Categories of uranium enrichment by weight fraction of 
        uranium-235.}
        \label{tab:enrichemnt}
        \begin{tabular}{l c c}
            \hline
            Category & Weight fraction (\%)\\\hline
            Depleted & $<$0.711 \\
            Natural & 0.711 \\
            LEU & 0.711-20 \\
            \gls{HALEU} & 5-20 \\
            \gls{HEU} & $\ge$20 \\
            \hline
        \end{tabular}
    \end{table}
    \end{columns}
\end{frame}

\begin{frame}
    \frametitle{Overview of the Nuclear Fuel Cycle}
    \begin{figure}[t]
    \centering
    \begin{tikzpicture}[node distance=1cm]
        \node (mine) [agent] {Mine};
        \node (mill) [agent, right of=mine] {Mill};
        \node (conversion) [agent, right of=mill, xshift=0.5cm] {Conversion};
        \node (enrichment) [agent, right of=conversion, xshift=1cm]{Enrichment};
        \node (fabrication) [agent, right of=enrichment, xshift=1.5cm]{Fuel Fabrication};
        \node (reactor) [agent, below of=fabrication, yshift=-1.5cm]{LWR};
        \node (storage) [agent,  left of=reactor, xshift=-1cm]{Storage};
        \node (sinkhlw) [agent, left of=storage, xshift=-1cm]{Repository};
        \node (sinkllw) [agent, above of=enrichment]{LLW Sink};


        \draw [arrow] (mine) --  (mill); 
        \draw [arrow] (mill) -- (conversion); 
        \draw [arrow] (conversion) -- (enrichment);
        \draw [arrow] (enrichment) -- (fabrication);
        \draw [arrow] (enrichment) -- (sinkllw);
        \draw [arrow] (fabrication) -- (reactor);
        \draw [arrow] (reactor) -- (storage);
        \draw [arrow] (storage) -- (sinkhlw);
        \pause
        \node (reprocessing) [transition, above of=storage]{Reprocessing};
        
        \draw [arrow] (storage) -- (reprocessing);
        \draw [arrow] (reprocessing) -- (fabrication);
        \draw [arrow] (reprocessing) -- (sinkhlw);
        \pause
        \node (downblending) [region, above of=fabrication]{Downblending};
        \draw [arrow] (downblending) -- (fabrication);

        \end{tikzpicture}
    \caption{Overview of the Nuclear Fuel Cycle.}
    \label{fig:fuel_cycle}
\end{figure}
\end{frame}

\begin{frame}
    \frametitle{Objectives}
    This work explores how developing a supply chain of \gls{HALEU} affects 
    the nuclear fuel cycle in the US. 
    Quantify material requirements of the transition to reactors 
              fueled by \gls{HALEU}
    \begin{itemize}
        \item<-1> The number of reactors deployed 
        \item<-1> Mass of enriched uranium sent to reactors
        \item<-1> Mass of feed uranium required to produce enriched uranium
        \item<-1> \gls{SWU} capacity required to enrich uranium 
    \end{itemize}
\end{frame}